\section{Divisibility in \(\mathbb{Z}\)}

\begin{definition}[Division]
  In a commutative ring \(R\), an element \(a \in R\) is said to divide an element \(b \in R\), written \(a \mid b\), if there exists an \(x \in R\) such that

\[
b=a x .
\]

We will sometimes write \(a \nmid b\) if \(a\) does not divide \(b\).
\end{definition}

\begin{remark}
  It is easy to show from the definition that if \(a \mid b\) and \(a \mid c\) for \(a, b, c \in R\), then \(a \mid(b x+c y)\), for any \(x, y \in R\). In particular, if \(a\) divides two elements, then it divides the sum and difference of those two elements.
\end{remark}

\begin{lemma}[Well-Ordering Principle (WOP)]
  Let \(S\) be a non-empty subset of \(\mathbb{N}_0=\mathbb{N} \cup\{0\}\). Then \(S\) has a smallest element.
\end{lemma}

\begin{proposition}
  Let \(a \in \mathbb{Z}, b \neq 0\). Then there exist in tegers \(q\) and \(r\) such that

\[
a=q \cdot b+r, \quad \text { and } \quad 0 \leq r<|b| .
\]

( \(q\) is called the quotient and \(r\) the remainder).
\end{proposition}

\begin{proof}
  Insert proof here
\end{proof}

\begin{example}
  Find the integer quotient and remainder of 435 divided by 43 in \(\mathbb{Z}\). Long division with remainder gives:

\[
435=43 \cdot 10+5,
\]

and \(5<43\), so we stop. So the integer quotient is \(q=10\) and the remainder is \(r=5\).
\end{example}

\begin{theorem}[Euclidean Algorithm in \(\mathbb{Z}\)]
  Let \(a, b \in \mathbb{Z}\), not both zero. Then \(\operatorname{gcd}(a, b)\) exists and can be computed by the Euclidean algorithm (given in the proof). Moreover, if \(d=\operatorname{gcd}(a, b)\), then there exist \(x, y \in \mathbb{Z}\) such that

\[
a x+b y=d .
\]
\end{theorem}

\begin{proof}
  Insert proof here
\end{proof}

\begin{example}
  
Example 2.5. We use the Euclidean Algorithm to compute \(d=\operatorname{gcd}(455,1235)\) and find integers \(x, y\) such that \(d=455 x+1235 y\).

Solution: Division with remainder:

\[
\begin{aligned}
1235 & =2 \cdot 455+325 \\
455 & =1 \cdot 325+130 \\
325 & =2 \cdot 130+65 \\
130 & =2 \cdot 65+0 .
\end{aligned}
\]

When the remainder is 0 , we stop. The gcd will be the last non-zero remainder, here \(d=65\). Solving for the remainders, substituting back and simplifying, we get

\[
\begin{aligned}
65 & =325-2 \cdot 130 \\
& =325-2(455-1 \cdot 325) \\
& =\cdots \\
& =455 \cdot(-8)+1235 \cdot 3 .
\end{aligned}
\]

so \(x=-8\) and \(y=3\).

\end{example}

\begin{lemma}[Primes that divide products divide at least one  factor]
  If \(p\) is a prime and \(p \mid\) ab for some \(a, b \in \mathbb{Z}\), then \(p \mid a\) or \(p \mid b\). (In other words, \(p\) prime \(\Rightarrow(\forall a, b \in \mathbb{Z}, p|a b \Rightarrow p| a\) or \(p \mid b)\).)
\end{lemma}

\begin{theorem}[Fundamental Theorem of Arithmetic]
Let \(n \in \mathbb{N}, n>1\). Then \(n\) can be written as a product of primes in a unique way (up to the order of the factors). More precisely, there are distinct primes \(p_1, \ldots, p_r\) and \(e_1, \ldots, e_r \in \mathbb{N}\) such that

\[
n=p_1^{e_1} \cdots p_r^{e_r},
\]

and if \(n=q_1^{f_1} \cdots q_s^{f_s}\) for some distinct primes \(q_1, \ldots, q_s\) and \(f_1, \ldots, f_s \in \mathbb{N}\), then \(r=s\), and up to re-ordering, \(p_i=q_i\) and \(e_i=f_i\) for all \(i\).
\end{theorem}

\section{Divisibility in \(F[x]\)}

\begin{definition}[Degree]
  For \(f(x)=a_0+a_1 x+\cdots+a_n x^n \in F[x]\), its \emph{degree} is

\[
\operatorname{deg} f:= \begin{cases}\max \left\{i \mid a_i \neq 0\right\} & \text { if } f(x) \neq 0, \\ -\infty & \text { if } f(x)=0 .\end{cases}
\]

Then we have the following basic properties for all \(f(x), g(x) \in F[x]\) :
\begin{itemize}
\item[-] \(\operatorname{deg}(f g)=\operatorname{deg} f+\operatorname{deg} g\),
\item[-] \(\operatorname{deg}(f+g) \leq \max \{\operatorname{deg} f, \operatorname{deg} g\}\). Note that if \(f(x) \neq g(x)\), then equality holds
\end{itemize}
\end{definition}

\begin{remark}
  We define \(\operatorname{deg}(f) = - \infty\) if \(f(x) = 0\) to preserve the properties of degree in the definition. 

If \(f=0\) and \(\operatorname{deg} 0=\operatorname{deg}(0 \cdot g)=\) \(\operatorname{deg} 0+\operatorname{deg} g\), then we can't allow \(\operatorname{deg} 0\) to be in \(\mathbb{Z}\), but it would work if \(\operatorname{deg} 0= \pm \infty\). However, if also \(\operatorname{deg} g=\operatorname{deg}(0+g) \leq \max \{\operatorname{deg} 0, \operatorname{deg} g\}\), for all \(g\), then we can't allow \(\operatorname{deg} 0=\infty\), but it does work if we set \(\operatorname{deg} 0=-\infty\).
\end{remark}


\begin{proposition}
Let \(f(x), g(x) \in F[x]\) with \(g(x) \neq 0\). Then there exist polynomials \(q(x), r(x) \in\) \(F[x]\) with \(\operatorname{deg} r<\operatorname{deg} g\) such that
\[
f(x)=q(x) g(x)+r(x) .
\]
\end{proposition}

\begin{example}[Long division of \(f(x)=x^3+x^2-3 x-3\) by \(g(x)=x^2+3 x+2\), with remainder]

  \begin{center}

\begin{tabular}[h]{|l|l|l|}
\hline term to get rid of & & \(x-2\) \\
\hline \(x^3\) & \(x^2+3 x+2\) & \[
\begin{aligned}
& x^3+x^2-3 x-3 \\
& -\left(x^3+3 x^2+2 x\right)
\end{aligned}
\] \\
\hline \(-2 x^2\) & & \(-2 x^2\)
\(-\left(-2 x^2-6 x-3\right.\)
\(-7)\) \\
\hline & & \(x+1\) \\
\hline
\end{tabular}

\end{center}

We have obtained a remainder of degree 1 , which is smaller than \(\operatorname{deg} g\), so we stop:

\[
f(x)=(x-2) g(x)+x+1 .
\]

Thus \(q(x)=x-2\) and \(r(x)=x+1\). Think of \(q(x)\) as the "quotient" and \(r(x)\) as the "remainder".
\end{example}

\begin{proof}
  insert proof here
\end{proof}


\section{GCD's and the Euclidean Algorithm in \(F[x]\)}


\begin{definition}[Monic Polynomial]
  A polynomial in \(F[x]\) is said to be monic if its leading coefficient is 1 .
\end{definition}

\begin{definition}[GCD for Polynomials]
Let \(f(x), g(x) \in F[x]\) with either \(f(x) \neq 0\) or \(g(x) \neq 0\). A greatest common divisor (GCD) of \(f(x)\) and \(g(x)\) is a polynomial of highest degree that divides both \(f(x)\) and \(g(x)\).
\end{definition}

\begin{example}
  Example 2.11. Let \(f(x)=x^2+1\) and \(g(x)=x^2+3 x+1\) in \(\mathbb{Q}[x]\). We perform the Euclidean algorithm to find a \(\operatorname{gcd}(f(x), g(x))\). Long division of \(f(x)\) by \(g(x)\) :


\begin{centering}

  \begin{tabular}[h!]{|l|l|l|}
\hline & 1 & \multirow{3}{*}{so \(f(x)=g(x)-3 x\)} \\
\hline \(x^2+3 x+1\) & \(x^2+1\) \(-\left(x^2+3 x+1\right)\) & \\
\hline & \(-3 x\) & \\
\hline
\end{tabular}

\end{centering}

Next, divide the divisor \(g(x)\) by the remainder \(-3 x\) :

\begin{centering}
  \begin{tabular}[h!]{|l|l|l|}
\hline & \(-\frac{1}{3} x+1\) & \multirow{3}{*}{so \(g(x)=\left(-\frac{1}{3} x+1\right)(-3 x)+1\)} \\
\hline \(-3 x\) & \[
\begin{gathered}
x^2+3 x+1 \\
-\left(x^2+1\right)
\end{gathered}
\] & \\
\hline & 1 & \\
\hline
\end{tabular}

\end{centering}


`
Next, divide the divisor \(-3 x\) from the previous step by the remainder 1: This obviously just gives

\[
-3 x=(-3 x) \times 1+0 .
\]

Now the remainder is 0 , so we stop. A GCD of \(f(x)\) and \(g(x)\) is given by the last non-zero remainder: 1.

Important for later: We can use this to write 1 as a linear combination of \(f(x)\) and \(g(x)\) :

\[
\begin{aligned}
1 & =g(x)-\left(-\frac{1}{3} x-1\right)(-3 x) \\
& =g(x)+\left(\frac{1}{3} x+1\right)(f(x)-g(x)) \\
& =\left(\frac{1}{3} x+1\right) f(x)-\frac{1}{3} x g(x) .
\end{aligned}
\]


\end{example}

\begin{definition}[Degree Function]
  Define the degree function \(d: R \backslash\{0\} \rightarrow \mathbb{N}\) by

\[
d(a):= \begin{cases}|a| & \text { if } a \in \mathbb{Z}, \\ \operatorname{deg} a & \text { if } a \in F[x] .\end{cases}
\]

For \(0 \in \mathbb{Z},|0|=0 \notin \mathbb{N}\) and for \(0 \in F[x], \operatorname{deg} 0=-\infty \notin \mathbb{N}\), which is why the domain of the function \(d\) is \(R \backslash\{0\}\), rather than \(R\)
\end{definition}

\begin{proposition}
  Let \(R\) be either \(\mathbb{Z}\) or \(F[x]\). Then, for any \(a, b \in R, b \neq 0\), there exist \(q, r \in R\) such that

\[
\begin{equation*}
a=q \cdot b+r, \quad r=0 \quad \text { or } \quad d(r)<d(b) \text {. } \tag{*}
\end{equation*}
\]
\end{proposition}

\begin{remark}
  The use of the degree function is to generalise the definition of the Euclidean algorithm for both polynomials and integer rings. This is to carefully define the inequalitiy in the definition.
\end{remark}

\begin{lemma}
  Let \(R\) be either \(\mathbb{Z}\) or \(F[x]\). If \(a, b \in R\) are such that \(a \mid b\), then \(d(a) \leq d(b)\).
\end{lemma}

\begin{proof}
  Suppose that \(a \mid b\), that is, \(b=a m\), for some \(m \in R\). Then \(d(b)=d(a m)=|a| \cdot|m| \geq|a|\) when \(R=\mathbb{Z}\) and \(d(b)=d(a m)=\operatorname{deg} a+\operatorname{deg}(m) \geq \operatorname{deg} a=d(a)\) when \(R=F[x]\).
\end{proof}

\begin{theorem}[Euclidean Algorithm - More Generally]
  
Let \(R\) be either \(\mathbb{Z}\) or \(F[x], F\) a field and let \(a, b \in R\) such that \(a \neq 0\) or \(b \neq 0\). Then
\begin{itemize}
\item[i)] A GCD of \(a\) and \(b\) exists and can be computed by the Euclidean algorithm (given in the proof);
\item[ii)] If \(R=F[x]\), the monic GCD is unique and we denote it by \(\operatorname{gcd}(a, b)\);
\item[iii)] If \(d\) is a GCD of \(a\) and \(b\), then there exist \(\alpha, \beta \in R\) such that \(\alpha a+\beta b=d\).
\end{itemize}

\end{theorem}

\begin{proof}
  insert proof here
\end{proof} 

\section{Units and Irreducibles, unique factorisation}

\begin{definition}[Invertible / Unit]
  Let \(R\) be a ring. An element \(a \in R\) is said to be invertible or a unit in \(R\) if there exists a \(b \in R\) such that \(a b=b a=1\) ( \(b\) is then the inverse of \(a\) and we write \(b=a^{-1}\) ). The set of units in \(R\) is denoted \(R^{\times}\).
\end{definition}

\begin{example}
  
  \begin{itemize}\emph{}


  \item[-] If \(a\) is a unit, then its inverse is unique: \(a b=1=a c=c a \Rightarrow c a b=c \Rightarrow b=c\).
\item[-] The units in \(\mathbb{Z}\) are \(\pm 1\) : if \(a \in \mathbb{Z}\) is a unit and \(a b=1\), then \(a \mid 1\).
\item[-] The units in \(F[x], F\) a field, are just the units in \(F\), that is, the non-zero constants. Proof: \(a(x) \in F[x]\) is a unit with \(a(x) b(x)=1\), for some \(b(x) \in F[x]\), then \(\operatorname{deg} a+\operatorname{deg} b=\operatorname{deg} 1=0\), and since \(\operatorname{deg} a, \operatorname{deg} b \geq 0\), we must have \(\operatorname{deg} a=\operatorname{deg} b=0\), which means that \(a\) (and \(b\) ) are non-zero constant polynomials.
\end{itemize}
\end{example}

\begin{definition}[Irreducible Polynomial]
A polynomial \(f(x) \in F[x]\) is irreducible if it is non-constant and does not have a proper factorisation, that is, \(f(x)\) cannot be written as \(f(x)=g(x) h(x)\) with \(\operatorname{deg} g>0\) and \(\operatorname{deg} h>0\). We generalise this definition to one that also captures the primes in \(\mathbb{Z}\).
\end{definition}

\begin{definition}[Irreducible Element of a Ring]
  Let \(R\) be a ring. An element \(r \in R\) is called irreducible if is satisfies
\begin{itemize}
\item[-] \(r\) is not a unit;
\item[-] If \(r=a b\), for some \(a, b \in R\), then \(a\) or \(b\) is a unit.
\end{itemize}
\end{definition}

\begin{lemma}[Generalised prime statement]
  Let \(R\) be either \(\mathbb{Z}\) or \(F[x], F\) a field. If \(p \in R\) is irreducible and \(p \mid\) ab for some \(a, b \in R\), then either \(p \mid a\) or \(p \mid b\).
\end{lemma}

\begin{proof}
  Insert proof here
\end{proof}

\begin{theorem}
 Let \(R\) be either \(\mathbb{Z}\) or \(F[x], F\) a field and let \(a \in R\) be a non-zero non-unit. Then a can be factorised uniquely into a product of irreducible elements, up to order of the factors and multiplication by invertible elements. 
\end{theorem}

\begin{proof}
  Insert proof here
\end{proof} 


